    \documentclass[../DoAn.tex]{subfiles}
    \begin{document}

    % Lệnh tạo khung ảnh trống — thay bằng \includegraphics khi đã có hình.
    \providecommand{\hinhcho}[1]{\fbox{\begin{minipage}[c][5.5cm][c]{0.85\textwidth}\centering\textit{[#1]}\end{minipage}}}

    Chương này trình bày quá trình phân tích thiết kế, xây dựng, kiểm thử và triển khai hệ thống dựa trên các yêu cầu đã xác định ở Chương \ref{chapter:Related_works} và nền tảng công nghệ đã lựa chọn ở Chương \ref{chapter:Methodology}. Nội dung lần lượt đi từ kiến trúc tổng thể, thiết kế chi tiết các thành phần, đến kết quả hiện thực hóa, kiểm thử và triển khai thực tế.

    \section{Thiết kế kiến trúc}

    \subsection{Lựa chọn kiến trúc phần mềm}

    Hệ thống được thiết kế theo mô hình phân tán, bao gồm ba thành phần chính: thiết bị nhúng (firmware ESP32 kết hợp module SIM7600), máy chủ xử lý (backend FastAPI) và ứng dụng di động (Flutter). Các thành phần được triển khai độc lập trên các môi trường khác nhau và giao tiếp thông qua mạng Internet.

    Để đảm bảo khả năng mở rộng và giảm sự phụ thuộc giữa các thành phần, hệ thống áp dụng mô hình hướng sự kiện (Event-driven Architecture). Thiết bị IoT gửi dữ liệu vị trí GPS và các bản tin cảnh báo đến MQTT Broker theo cơ chế publish/subscribe. Backend đăng ký nhận các bản tin từ Broker để xử lý nghiệp vụ, lưu trữ dữ liệu và cung cấp dịch vụ cho ứng dụng di động. Đối với giao tiếp giữa backend và ứng dụng di động, hệ thống sử dụng kết hợp REST API và WebSocket. REST API được sử dụng để thực hiện các thao tác như đăng nhập, quản lý tài khoản, thiết bị và truy vấn dữ liệu, trong khi WebSocket được sử dụng để truyền dữ liệu vị trí theo thời gian thực từ backend đến ứng dụng mà không cần gửi yêu cầu liên tục.

    Việc sử dụng MQTT kết hợp với WebSocket giúp tách biệt hoàn toàn giữa thiết bị, backend và ứng dụng di động. Thiết bị chỉ cần gửi dữ liệu đến MQTT Broker mà không cần biết trạng thái hoạt động của backend hoặc ứng dụng. Backend sẽ xử lý dữ liệu ngay khi nhận được bản tin từ Broker và chủ động đẩy dữ liệu mới đến ứng dụng thông qua WebSocket. Kiến trúc này giúp giảm lưu lượng trao đổi dữ liệu, nâng cao khả năng mở rộng và đáp ứng tốt yêu cầu cập nhật vị trí theo thời gian thực.

    Bên trong backend, hệ thống được tổ chức theo kiến trúc phân tầng (Layered Architecture), trong đó mỗi tầng đảm nhiệm một chức năng riêng biệt nhằm tăng tính mô-đun và thuận tiện cho việc bảo trì. Cụ thể:

    \begin{itemize}
        \item \textbf{Tầng API} (gói \texttt{api}): tiếp nhận các yêu cầu HTTP và WebSocket từ ứng dụng di động, kiểm tra xác thực bằng JWT và chuyển tiếp yêu cầu đến tầng dịch vụ.
        \item \textbf{Tầng dịch vụ} (gói \texttt{services}): chứa các nghiệp vụ chính của hệ thống như xử lý dữ liệu GPS, kiểm tra geofence, quản lý thiết bị, xử lý bản tin MQTT và gửi cảnh báo.
        \item \textbf{Tầng truy xuất dữ liệu} (gói \texttt{repositories}): thực hiện các thao tác đọc và ghi dữ liệu với cơ sở dữ liệu MongoDB.
        \item \textbf{Tầng dữ liệu và hạ tầng} (gói \texttt{models}, \texttt{db}, \texttt{core}, \texttt{ws}): định nghĩa mô hình dữ liệu, quản lý kết nối cơ sở dữ liệu, cấu hình hệ thống, xác thực JWT và quản lý các kết nối WebSocket.
    \end{itemize}

    Ứng dụng di động được xây dựng bằng Flutter và tổ chức theo hướng tách biệt giữa giao diện (\texttt{screens}, \texttt{widgets}), quản lý trạng thái (\texttt{providers}), dịch vụ giao tiếp với backend (\texttt{services}) và mô hình dữ liệu (\texttt{models}). Cách tổ chức này giúp mã nguồn rõ ràng, thuận tiện cho việc phát triển, bảo trì và mở rộng ứng dụng trong tương lai.

    Hình \ref{fig:kientruc_tongthe} minh họa kiến trúc tổng thể của hệ thống, bao gồm các thành phần chính và luồng trao đổi dữ liệu giữa thiết bị IoT, MQTT Broker, backend và ứng dụng di động.

    \begin{figure}[H]
        \centering
        \includegraphics[width=0.95\textwidth]{Hinhve/kientruc_tongthe.png}
        \caption{Kiến trúc tổng thể của hệ thống theo các tầng chức năng}
        \label{fig:kientruc_tongthe}
    \end{figure}

    \subsection{Thiết kế tổng quan}

    Hình \ref{fig:bieudo_goi} trình bày biểu đồ gói (Package Diagram) của toàn bộ hệ thống, bao gồm ba thành phần chính là firmware ESP32, backend FastAPI và ứng dụng di động Flutter, cùng với các thành phần bên ngoài như MQTT Broker, MongoDB và Telegram Bot. Các gói được tổ chức theo chức năng và phân thành nhiều tầng nhằm giảm sự phụ thuộc giữa các thành phần, giúp mã nguồn dễ bảo trì, thuận tiện cho việc phát triển và mở rộng hệ thống.

    \begin{figure}[H]
        \centering
        \includegraphics[width=0.95\textwidth]{Hinhve/bieudo_goi.png}
        \caption{Biểu đồ gói tổng quan của hệ thống}
        \label{fig:bieudo_goi}
    \end{figure}

    Nhiệm vụ của từng nhóm gói được tóm tắt như sau:

    \textbf{Backend}
    \begin{itemize}
        \item \texttt{api}: tiếp nhận các yêu cầu REST API và WebSocket từ ứng dụng di động, kiểm tra xác thực JWT và chuyển tiếp yêu cầu đến tầng xử lý nghiệp vụ.
        \item \texttt{services}: xử lý các nghiệp vụ của hệ thống như tiếp nhận dữ liệu GPS, kiểm tra geofence, quản lý thiết bị, xử lý bản tin MQTT và gửi cảnh báo.
        \item \texttt{repositories}: thực hiện các thao tác đọc và ghi dữ liệu với cơ sở dữ liệu MongoDB.
        \item \texttt{models}: định nghĩa các mô hình dữ liệu được sử dụng trong hệ thống.
        \item \texttt{db}: quản lý kết nối với cơ sở dữ liệu MongoDB.
        \item \texttt{core}: quản lý cấu hình hệ thống, xác thực JWT và các chức năng dùng chung.
        \item \texttt{ws}: quản lý các kết nối WebSocket phục vụ việc cập nhật dữ liệu theo thời gian thực.
    \end{itemize}

    \textbf{Flutter App}
    \begin{itemize}
        \item \texttt{screens}: xây dựng các màn hình chức năng của ứng dụng.
        \item \texttt{widgets}: chứa các thành phần giao diện có thể tái sử dụng.
        \item \texttt{providers}: quản lý trạng thái của ứng dụng.
        \item \texttt{services}: thực hiện giao tiếp với backend thông qua REST API và WebSocket.
        \item \texttt{models}: định nghĩa các mô hình dữ liệu.
        \item \texttt{config} và \texttt{utils}: quản lý cấu hình và các hàm tiện ích của ứng dụng.
    \end{itemize}

    \textbf{ESP32 Firmware}
    \begin{itemize}
        \item \texttt{gps}: thu thập dữ liệu vị trí từ module SIM7600.
        \item \texttt{mqtt\_client}: thiết lập kết nối với MQTT Broker và gửi dữ liệu vị trí, tín hiệu SOS của thiết bị lên hệ thống.
    \end{itemize}

    Quan hệ giữa các thành phần phản ánh đúng kiến trúc của hệ thống. Firmware ESP32 thực hiện publish dữ liệu vị trí và tín hiệu SOS lên MQTT Broker. Backend đăng ký (subscribe) các topic tương ứng để tiếp nhận và xử lý dữ liệu. Sau khi xử lý, backend lưu thông tin vào MongoDB, gửi cảnh báo đến Telegram Bot khi phát hiện các sự kiện bất thường, đồng thời truyền dữ liệu vị trí theo thời gian thực đến ứng dụng di động thông qua WebSocket. Bên cạnh đó, ứng dụng Flutter sử dụng các REST API để thực hiện các chức năng như đăng nhập, quản lý thiết bị, thiết lập vùng an toàn, truy vấn lịch sử vị trí và các chức năng quản lý khác.

    Biểu đồ gói cho thấy sự phân tách rõ ràng giữa các thành phần của hệ thống cũng như mối quan hệ giữa chúng. Việc tổ chức các gói theo chức năng giúp giảm sự phụ thuộc giữa các thành phần, tăng tính mô-đun và tạo thuận lợi cho quá trình bảo trì, kiểm thử cũng như mở rộng hệ thống trong tương lai.

    \subsection{Thiết kế chi tiết gói và thiết bị nhúng}

    Trong backend, phần lớn các chức năng nghiệp vụ được triển khai tại gói \texttt{services}. Đây là thành phần trung tâm của hệ thống, chịu trách nhiệm xử lý dữ liệu nhận từ thiết bị IoT, kiểm tra điều kiện cảnh báo, quản lý vùng an toàn và điều phối việc lưu trữ dữ liệu cũng như gửi thông báo. Các dịch vụ trong gói này phối hợp với tầng \texttt{repositories} để truy xuất cơ sở dữ liệu MongoDB và với thành phần quản lý WebSocket để truyền dữ liệu vị trí theo thời gian thực đến ứng dụng di động.

    Về phía thiết bị nhúng, firmware được xây dựng trên nền tảng ESP32 và tổ chức thành các mô-đun chức năng riêng biệt nhằm thực hiện việc thu thập dữ liệu GPS, quản lý kết nối mạng di động và truyền dữ liệu đến hệ thống thông qua giao thức MQTT. Cấu trúc này giúp việc phát triển và bảo trì firmware trở nên thuận tiện, đồng thời dễ dàng mở rộng khi bổ sung thêm các cảm biến hoặc chức năng mới.

    Hình \ref{fig:phancung} trình bày sơ đồ kết nối phần cứng của thiết bị. ESP32 đóng vai trò là bộ điều khiển trung tâm, giao tiếp với module SIM7600 thông qua giao tiếp UART để thu thập dữ liệu GPS và truyền dữ liệu qua mạng di động. Ngoài ra, hệ thống sử dụng một nút nhấn SOS kết nối với chân GPIO của ESP32 để người dùng có thể gửi tín hiệu khẩn cấp. Module SIM7600 được kết nối với anten GPS và anten GSM nhằm đảm bảo khả năng thu tín hiệu vệ tinh và kết nối mạng di động ổn định.

    \begin{figure}[H]
        \centering
        \includegraphics[width=0.95\textwidth]{Hinhve/phancung.png}
        \caption{Sơ đồ kết nối phần cứng của thiết bị}
        \label{fig:phancung}
    \end{figure}

    Hình \ref{fig:firmware_setup} mô tả quá trình khởi tạo của firmware. Sau khi cấp nguồn, ESP32 tiến hành khởi tạo các ngoại vi cần thiết, thiết lập kết nối với module SIM7600, kích hoạt chức năng GPS, kết nối mạng di động và thiết lập kết nối với MQTT Broker. Sau khi các bước khởi tạo hoàn tất, hệ thống bắt đầu chuyển sang chế độ hoạt động bình thường.

    \begin{figure}[H]
        \centering
        \includegraphics[width=0.72\textwidth]{Hinhve/firmware_setup.png}
        \caption{Lưu đồ giai đoạn khởi tạo (Setup) của firmware ESP32}
        \label{fig:firmware_setup}
    \end{figure}

    Hình \ref{fig:firmware_loop} mô tả vòng lặp chính của firmware. Trong quá trình hoạt động, thiết bị duy trì kết nối với MQTT Broker, định kỳ đọc dữ liệu vị trí từ module GPS và gửi dữ liệu lên backend thông qua MQTT khi thu được tọa độ hợp lệ. Đồng thời, firmware liên tục theo dõi trạng thái của nút SOS để kịp thời xử lý khi xảy ra tình huống khẩn cấp.

    \begin{figure}[H]
        \centering
        \includegraphics[width=0.72\textwidth]{Hinhve/firmware_loop.png}
        \caption{Lưu đồ vòng lặp chính (Main Loop) của firmware ESP32}
        \label{fig:firmware_loop}
    \end{figure}

    Hình \ref{fig:firmware_sos} mô tả quy trình xử lý sự kiện SOS. Khi người dùng nhấn nút SOS, firmware ghi nhận sự kiện, thu thập thông tin vị trí hiện tại và gửi bản tin cảnh báo đến backend thông qua MQTT. Sau khi backend tiếp nhận bản tin, hệ thống thực hiện lưu trữ dữ liệu và gửi thông báo qua Telegram Bot để phụ huynh có thể nhận được cảnh báo trong thời gian sớm nhất.

    \begin{figure}[H]
        \centering
        \includegraphics[width=0.72\textwidth]{Hinhve/firmware_sos.png}
        \caption{Lưu đồ xử lý ngắt SOS của firmware ESP32}
        \label{fig:firmware_sos}
    \end{figure}

    \section{Thiết kế chi tiết}

    \subsection{Thiết kế giao diện}

    Ứng dụng di động được phát triển bằng Flutter và hỗ trợ cả hai nền tảng Android và iOS. Nhờ sử dụng cơ chế bố cục co giãn (responsive layout), giao diện có thể thích ứng với nhiều kích thước màn hình khác nhau, đồng thời đảm bảo tính thống nhất về cách hiển thị và trải nghiệm người dùng trên các thiết bị.

    Trong quá trình thiết kế giao diện, hệ thống áp dụng một số nguyên tắc nhằm đảm bảo tính trực quan và dễ sử dụng. Màu sắc được sử dụng thống nhất giữa các màn hình, trong đó màu chủ đạo được áp dụng cho thanh điều hướng và các nút chức năng chính, còn màu cảnh báo được sử dụng để biểu thị các sự kiện như vi phạm vùng an toàn hoặc tín hiệu SOS. Các thành phần điều khiển được bố trí phù hợp với thao tác trên màn hình cảm ứng, giúp người dùng dễ dàng thực hiện các chức năng theo dõi và quản lý thiết bị.

    Ứng dụng bao gồm bốn màn hình chính: màn hình \textbf{xác thực} (đăng nhập, đăng ký tài khoản), màn hình \textbf{bảng điều khiển} (theo dõi bản đồ, quản lý thiết bị, thiết lập vùng an toàn), màn hình \textbf{thống kê} (xem lịch sử di chuyển) và màn hình \textbf{hồ sơ} (thông tin người dùng). Trong đó, màn hình bảng điều khiển là giao diện trung tâm của hệ thống, cho phép phụ huynh theo dõi vị trí hiện tại của trẻ trên bản đồ, xem trạng thái thiết bị, nhận các cảnh báo theo thời gian thực và truy cập nhanh đến các chức năng quản lý khác.

    Hình \ref{fig:tk_giaodien} minh họa thiết kế giao diện của màn hình bảng điều khiển. Giao diện được bố trí theo hướng trực quan, trong đó khu vực bản đồ chiếm phần lớn diện tích màn hình để hiển thị vị trí hiện tại của thiết bị và các vùng an toàn đã được thiết lập. Các khu vực chức năng được sắp xếp xung quanh bản đồ, bao gồm danh sách thiết bị, thông tin trạng thái, các nút điều khiển và các chức năng quản lý. Khi backend phát hiện các sự kiện như cập nhật vị trí hoặc vi phạm geofence, thông tin sẽ được gửi đến ứng dụng thông qua WebSocket và hiển thị ngay trên giao diện; riêng tín hiệu SOS được gửi ngay đến phụ huynh qua Telegram Bot. Thiết kế này giúp người dùng có thể theo dõi trạng thái của thiết bị theo thời gian thực, đồng thời thao tác nhanh chóng với các chức năng quản lý ngay trên cùng một màn hình.

    \begin{figure}[H]
        \centering
        \includegraphics[width=0.85\textwidth]{Hinhve/tk_giaodien.png}
        \caption{Thiết kế giao diện màn hình bảng điều khiển}
        \label{fig:tk_giaodien}
    \end{figure}

    \subsection{Thiết kế lớp}

    Phần này trình bày thiết kế các lớp chính tham gia xử lý nghiệp vụ trong backend. Mỗi lớp đảm nhận một nhóm chức năng riêng và phối hợp với các lớp khác nhằm thực hiện quá trình tiếp nhận dữ liệu từ thiết bị IoT, xử lý nghiệp vụ, lưu trữ dữ liệu và cập nhật thông tin theo thời gian thực đến ứng dụng di động.

    \textbf{Lớp \texttt{LocationProcessor}} đóng vai trò là trung tâm điều phối quá trình xử lý dữ liệu vị trí. Lớp này tiếp nhận dữ liệu GPS nhận được từ MQTT Broker, kiểm tra tính hợp lệ của dữ liệu, điều phối việc kiểm tra vùng an toàn, lưu thông tin vị trí vào cơ sở dữ liệu và gửi dữ liệu thời gian thực tới các ứng dụng đang kết nối. Trong quá trình xử lý, lớp phối hợp với \texttt{GeofenceService}, \texttt{AlertStateService}, \texttt{LocationRepository} và \texttt{ConnectionManager} để hoàn thành toàn bộ quy trình xử lý bản tin vị trí.

    \textbf{Lớp \texttt{GeofenceService}} chịu trách nhiệm quản lý các vùng an toàn và thực hiện các phép tính xác định vị trí của thiết bị so với vùng giám sát. Lớp này hỗ trợ các chế độ geofence của hệ thống và sử dụng thuật toán Haversine để tính khoảng cách giữa các tọa độ địa lý trong quá trình kiểm tra.

    \textbf{Lớp \texttt{AlertStateService}} quản lý trạng thái cảnh báo của từng thiết bị và quyết định thời điểm phát sinh hoặc kết thúc cảnh báo. Thành phần này giúp hạn chế việc gửi lặp các cảnh báo do sai số GPS hoặc do thiết bị liên tục dao động quanh ranh giới của vùng an toàn, đồng thời theo dõi trạng thái cảnh báo để đảm bảo mỗi sự kiện chỉ được xử lý khi cần thiết.

    \textbf{Lớp \texttt{ConnectionManager}} thuộc gói \texttt{ws} chịu trách nhiệm quản lý các kết nối WebSocket giữa backend và ứng dụng di động. Khi có dữ liệu vị trí mới hoặc phát sinh cảnh báo, lớp này sẽ gửi thông tin tới các ứng dụng đang theo dõi thiết bị tương ứng, giúp dữ liệu được cập nhật theo thời gian thực.

    Ngoài các lớp xử lý nghiệp vụ, backend còn sử dụng các lớp thuộc gói \texttt{repositories} để thực hiện việc truy xuất dữ liệu với MongoDB và các lớp trong gói \texttt{models} để định nghĩa cấu trúc dữ liệu được trao đổi giữa các thành phần của hệ thống. Việc tách biệt các lớp theo chức năng giúp giảm sự phụ thuộc giữa các thành phần, tăng khả năng tái sử dụng mã nguồn và thuận tiện cho việc bảo trì cũng như mở rộng hệ thống trong tương lai.

    \subsection{Biểu đồ tuần tự}

    Biểu đồ tuần tự (Sequence Diagram) được sử dụng để mô tả trình tự trao đổi thông điệp giữa các thành phần trong hệ thống đối với các kịch bản nghiệp vụ quan trọng. Thông qua các biểu đồ này có thể thấy rõ vai trò của từng thành phần, thứ tự thực hiện các thao tác cũng như luồng dữ liệu từ thiết bị IoT đến backend và ứng dụng di động.

    Hình \ref{fig:sd01} mô tả quy trình cập nhật vị trí theo thời gian thực (SD01). Thiết bị ESP32 thu thập dữ liệu GPS và gửi lên MQTT Broker. Backend đăng ký (subscribe) các topic tương ứng để tiếp nhận bản tin, thực hiện kiểm tra dữ liệu, lưu thông tin vị trí vào MongoDB và gửi dữ liệu mới đến ứng dụng Flutter thông qua WebSocket. Nhờ sử dụng WebSocket, vị trí của thiết bị được cập nhật liên tục trên bản đồ mà không cần ứng dụng gửi yêu cầu truy vấn định kỳ.

    \begin{figure}[H]
        \centering
        \includegraphics[width=0.95\textwidth]{Hinhve/sequence_vitri.png}
        \caption{SD01 – Biểu đồ tuần tự: Cập nhật vị trí GPS thời gian thực}
        \label{fig:sd01}
    \end{figure}

    Hình \ref{fig:sd02} mô tả quy trình xử lý khi thiết bị vi phạm vùng an toàn (SD02). Sau khi nhận được dữ liệu vị trí mới, backend thực hiện kiểm tra geofence. Nếu phát hiện thiết bị đã ra khỏi vùng an toàn, hệ thống tạo bản ghi cảnh báo, lưu vào cơ sở dữ liệu, gửi thông báo đến Telegram Bot và đồng thời cập nhật trạng thái cảnh báo tới ứng dụng Flutter thông qua WebSocket. Quy trình này giúp phụ huynh nhận được cảnh báo gần như ngay sau khi sự kiện xảy ra.

    \begin{figure}[H]
        \centering
        \includegraphics[width=0.95\textwidth]{Hinhve/sequence_geofence.png}
        \caption{SD02 – Biểu đồ tuần tự: Cảnh báo vi phạm vùng an toàn}
        \label{fig:sd02}
    \end{figure}

    Hình \ref{fig:sd03} mô tả quy trình xử lý tín hiệu SOS khẩn cấp (SD03). Khi người dùng nhấn nút SOS trên thiết bị, firmware gửi bản tin khẩn cấp đến MQTT Broker. Backend tiếp nhận bản tin, lưu thông tin cảnh báo và gửi thông báo khẩn cấp kèm vị trí hiện tại của thiết bị đến Telegram Bot để phụ huynh nhanh chóng nắm bắt tình huống.

    \begin{figure}[H]
        \centering
        \includegraphics[width=0.95\textwidth]{Hinhve/sequence_sos.png}
        \caption{SD03 – Biểu đồ tuần tự: Gửi tín hiệu SOS khẩn cấp}
        \label{fig:sd03}
    \end{figure}

    Ba biểu đồ tuần tự trên thể hiện ba luồng xử lý quan trọng nhất của hệ thống, bao gồm cập nhật vị trí, phát hiện vi phạm vùng an toàn và xử lý tín hiệu SOS. Các biểu đồ cho thấy sự phối hợp giữa thiết bị nhúng, MQTT Broker, backend, MongoDB, Telegram Bot và ứng dụng Flutter trong việc xử lý dữ liệu và cung cấp thông tin theo thời gian thực.

    \subsection{Thiết kế cơ sở dữ liệu}

    Hệ thống sử dụng MongoDB làm hệ quản trị cơ sở dữ liệu theo mô hình NoSQL hướng tài liệu. Mặc dù MongoDB không yêu cầu lược đồ dữ liệu cố định, đề tài vẫn xây dựng biểu đồ thực thể liên kết (E-R) nhằm mô tả cấu trúc logic của dữ liệu và mối liên hệ giữa các collection trong hệ thống. Biểu đồ được trình bày ở Hình \ref{fig:erd}.

    \begin{figure}[H]
        \centering
        \includegraphics[width=0.95\textwidth]{Hinhve/er_diagram.png}
        \caption{Biểu đồ thực thể liên kết (E-R) của cơ sở dữ liệu}
        \label{fig:erd}
    \end{figure}

    Cơ sở dữ liệu gồm sáu collection chính với chức năng như sau:

    \begin{itemize}
        \item \texttt{users}: lưu thông tin tài khoản của phụ huynh, bao gồm địa chỉ email, mật khẩu đã được băm và thời điểm tạo tài khoản.
        \item \texttt{device\_permissions}: quản lý quyền truy cập thiết bị của người dùng. Mỗi bản ghi liên kết một người dùng với một thiết bị và xác định quyền tương ứng như chủ sở hữu (owner) hoặc người được chia sẻ (shared). Một người dùng có thể được cấp quyền trên nhiều thiết bị và một thiết bị cũng có thể được chia sẻ cho nhiều người dùng.
        \item \texttt{device\_configs}: lưu các thông tin cấu hình của thiết bị như email phụ huynh, trạng thái bật hoặc tắt cảnh báo và các tham số cấu hình khác phục vụ quá trình vận hành.
        \item \texttt{locations}: lưu lịch sử vị trí GPS của thiết bị theo thời gian, bao gồm tọa độ, thời gian ghi nhận, trạng thái geofence và các thông tin liên quan phục vụ việc truy vết và thống kê hành trình.
        \item \texttt{geofence\_configs}: lưu cấu hình các vùng an toàn của từng thiết bị, bao gồm các thông tin như chế độ hoạt động, tọa độ tâm, bán kính hoặc danh sách các điểm tạo thành tuyến đường.
        \item \texttt{alerts}: lưu lịch sử các cảnh báo được phát sinh trong quá trình hoạt động của hệ thống, bao gồm cảnh báo vi phạm vùng an toàn và tín hiệu SOS cùng với thời gian và vị trí xảy ra sự kiện.
    \end{itemize}

    Để nâng cao hiệu năng truy vấn, collection \texttt{locations} sử dụng chỉ mục kết hợp trên các trường \texttt{deviceId} và \texttt{timestamp}, giúp tăng tốc việc truy vấn lịch sử vị trí của từng thiết bị theo khoảng thời gian. Ngoài ra, hệ thống sử dụng chỉ mục TTL trên trường \texttt{receivedAt} để tự động xóa các bản ghi đã quá thời gian lưu trữ theo quy định, giúp giảm dung lượng lưu trữ và duy trì hiệu năng của cơ sở dữ liệu. Collection \texttt{users} được tạo chỉ mục duy nhất trên trường \texttt{email} nhằm đảm bảo mỗi tài khoản sử dụng một địa chỉ email duy nhất.

    \section{Xây dựng ứng dụng}

    \subsection{Thư viện và công cụ sử dụng}

    Bảng \ref{table:congcu} liệt kê các công cụ, ngôn ngữ và thư viện chính được sử dụng để phát triển ba thành phần của hệ thống cùng phiên bản tương ứng.

    \begin{table}[H]
    \centering
    \caption{Danh sách thư viện và công cụ sử dụng}
    \label{table:congcu}
    \begin{tabular}{|l|l|l|}
    \hline
    \textbf{Thành phần} & \textbf{Công cụ / Thư viện} & \textbf{Phiên bản} \\ \hline
    Backend & Python & 3.12 \\ \hline
    Backend & FastAPI & 0.115.12 \\ \hline
    Backend & Uvicorn & 0.34.2 \\ \hline
    Backend & paho-mqtt & 2.1.0 \\ \hline
    Backend & PyMongo & 4.12.0 \\ \hline
    Backend & Pydantic & 2.12.0 \\ \hline
    Backend & python-jose (JWT) & 3.3.0 \\ \hline
    Backend & passlib/bcrypt & 1.7.4 / 4.0.1 \\ \hline
    Backend & pycryptodome (AES) & 3.21.0 \\ \hline
    Ứng dụng & Flutter / Dart SDK & Dart $\geq$ 3.0 \\ \hline
    Ứng dụng & flutter\_map & 6.1.0 \\ \hline
    Ứng dụng & latlong2 & 0.9.0 \\ \hline
    Ứng dụng & web\_socket\_channel & 2.4.0 \\ \hline
    Ứng dụng & provider & 6.1.1 \\ \hline
    Ứng dụng & geolocator & 10.1.0 \\ \hline
    Ứng dụng & fl\_chart & 0.68.0 \\ \hline
    Ứng dụng & http & 1.1.0 \\ \hline
    Firmware & Arduino framework (ESP32) & --- \\ \hline
    Firmware & PubSubClient & 2.8 \\ \hline
    Firmware & TinyGSM & 0.11.7 \\ \hline
    Hạ tầng & MongoDB / Atlas & 7.x \\ \hline
    Hạ tầng & EMQX Public Broker & --- \\ \hline
    Công cụ & VS Code, PlatformIO, Git & --- \\ \hline
    \end{tabular}
    \end{table}

    \subsection{Kết quả xây dựng}

    Đề tài đã xây dựng hoàn chỉnh hệ thống theo mô hình gồm ba thành phần chính: thiết bị nhúng, máy chủ backend và ứng dụng di động. Các thành phần được phát triển độc lập nhưng có khả năng phối hợp với nhau thông qua giao thức MQTT, REST API và WebSocket để tạo thành một hệ thống theo dõi vị trí hoạt động hoàn chỉnh.

    \begin{itemize}
        \item \textbf{Firmware thiết bị:} được triển khai trên ESP32 kết hợp với module SIM7600, thực hiện đọc dữ liệu GPS, xử lý tín hiệu từ nút SOS và truyền dữ liệu vị trí cũng như cảnh báo khẩn cấp lên MQTT Broker thông qua mạng di động 4G.
        \item \textbf{Backend:} được xây dựng bằng FastAPI, đảm nhiệm việc tiếp nhận dữ liệu từ MQTT Broker, xử lý nghiệp vụ, lưu trữ dữ liệu vào MongoDB, cung cấp REST API cho ứng dụng di động và gửi dữ liệu thời gian thực thông qua WebSocket. Ngoài ra, backend còn tích hợp Telegram Bot để gửi cảnh báo khi phát hiện vi phạm vùng an toàn hoặc nhận tín hiệu SOS.
        \item \textbf{Ứng dụng di động:} được phát triển bằng Flutter, hỗ trợ Android và iOS, cho phép người dùng đăng nhập, theo dõi vị trí thiết bị trên bản đồ, cấu hình vùng an toàn, xem lịch sử di chuyển và nhận cảnh báo theo thời gian thực.
    \end{itemize}

    Để đánh giá quy mô của hệ thống, Bảng \ref{table:thongke} thống kê số lượng tệp mã nguồn và số dòng mã của từng thành phần.

    \begin{table}[H]
    \centering
    \caption{Thống kê quy mô mã nguồn}
    \label{table:thongke}
    \begin{tabular}{|l|c|c|}
    \hline
    \textbf{Thành phần} & \textbf{Số tệp} & \textbf{Số dòng mã (xấp xỉ)} \\ \hline
    Backend (Python)        & 31 & 2.337 \\ \hline
    Ứng dụng (Dart)         & 27 & 5.348 \\ \hline
    Firmware (C++/Arduino)  & 3  & 570 \\ \hline
    \textbf{Tổng}           & \textbf{61} & \textbf{8.255} \\ \hline
    \end{tabular}
    \end{table}

    Tổng cộng hệ thống gồm khoảng 61 tệp mã nguồn với hơn 8.000 dòng mã. Trong đó, backend được tổ chức thành các gói \texttt{api}, \texttt{services}, \texttt{repositories}, \texttt{core}, \texttt{ws}, \texttt{models} và \texttt{db}; ứng dụng Flutter được chia thành các nhóm \texttt{screens}, \texttt{providers}, \texttt{services}, \texttt{widgets} và \texttt{models}; firmware được xây dựng theo các mô-đun xử lý GPS và truyền thông MQTT. Cách tổ chức này giúp mã nguồn có cấu trúc rõ ràng, thuận tiện cho việc bảo trì và mở rộng trong tương lai.

    \subsection{Minh họa các chức năng chính}

    Phần này minh họa các chức năng quan trọng nhất của hệ thống qua ảnh chụp màn hình thực tế.

    \begin{figure}[H]
        \centering
        \includegraphics[width=0.75\textwidth]{Hinhve/demo_auth.png}
        \caption{Màn hình đăng nhập và đăng ký tài khoản}
        \label{fig:demo_auth}
    \end{figure}

    \begin{figure}[H]
        \centering
        \includegraphics[width=0.75\textwidth]{Hinhve/demo_map.png}
        \caption{Theo dõi vị trí trẻ theo thời gian thực trên bản đồ}
        \label{fig:demo_map}
    \end{figure}

    \begin{figure}[H]
        \centering
        \includegraphics[width=0.75\textwidth]{Hinhve/demo_geofence.png}
        \caption{Tạo và chỉnh sửa vùng an toàn trên bản đồ}
        \label{fig:demo_geofence}
    \end{figure}

    \begin{figure}[H]
        \centering
        \includegraphics[width=0.75\textwidth]{Hinhve/demo_alert.png}
        \caption{Thông báo cảnh báo vi phạm vùng an toàn và tín hiệu SOS}
        \label{fig:demo_alert}
    \end{figure}

    \begin{figure}[H]
        \centering
        \includegraphics[width=0.75\textwidth]{Hinhve/demo_stats.png}
        \caption{Thống kê hành trình và số lần vi phạm}
        \label{fig:demo_stats}
    \end{figure}

    \section{Kiểm thử}

    Hệ thống được kiểm thử chủ yếu bằng kỹ thuật \textbf{kiểm thử hộp đen} (black-box testing) ở mức chức năng và \textbf{kiểm thử tích hợp} đầu cuối, do tính chất hệ thống phân tán gồm nhiều thành phần giao tiếp qua mạng. Việc kiểm thử được thực hiện thủ công kết hợp với một thiết bị mô phỏng publish bản tin MQTT để tái lập các kịch bản di chuyển mà không cần thiết bị vật lý.

    Đề tài thiết kế các trường hợp kiểm thử cho ba chức năng quan trọng nhất: cập nhật vị trí thời gian thực, cảnh báo vi phạm geofence, và gửi tín hiệu SOS. Bảng \ref{table:kiemthu} trình bày một số trường hợp tiêu biểu.

    \begin{table}[H]
    \centering
    \caption{Một số trường hợp kiểm thử tiêu biểu}
    \label{table:kiemthu}
    \begin{tabular}{|c|p{4.2cm}|p{4.2cm}|c|}
    \hline
    \textbf{Mã} & \textbf{Mô tả} & \textbf{Kết quả mong đợi} & \textbf{KQ} \\ \hline
    TC01 & Thiết bị gửi tọa độ hợp lệ & Bản đồ cập nhật vị trí trong dưới 3 giây & Đạt \\ \hline
    TC02 & Thiết bị mất tín hiệu GPS & Không publish, giữ vị trí cuối & Đạt \\ \hline
    TC03 & Trẻ ra ngoài vùng an toàn & Sinh cảnh báo, gửi Telegram & Đạt \\ \hline
    TC04 & Trẻ quay lại vùng an toàn & Gửi thông báo trở về vùng an toàn & Đạt \\ \hline
    TC05 & Hai cảnh báo liên tiếp trong cooldown & Chỉ gửi một cảnh báo & Đạt \\ \hline
    TC06 & Dao động GPS nhỏ quanh ranh giới & Bỏ qua nhiễu, không cảnh báo sai & Đạt \\ \hline
    TC07 & Nhấn nút SOS & Gửi cảnh báo SOS kèm link bản đồ qua Telegram & Đạt \\ \hline
    TC08 & Đăng nhập sai mật khẩu & Báo lỗi, không cấp token & Đạt \\ \hline
    TC09 & Truy cập API không kèm token hợp lệ & Từ chối truy cập (401) & Đạt \\ \hline
    TC10 & Mất kết nối WebSocket & Ứng dụng tự kết nối lại & Đạt \\ \hline
    \end{tabular}
    \end{table}

    Tổng cộng 10 trường hợp kiểm thử cho ba chức năng cốt lõi đều cho kết quả đạt, cho thấy hệ thống đáp ứng đúng các yêu cầu chức năng đã đặc tả ở Chương \ref{chapter:Related_works}. Các trường hợp kiểm thử chi tiết khác được trình bày trong phần phụ lục.

    \section{Triển khai}

    Hệ thống được triển khai theo mô hình đám mây cho phần backend và sử dụng các dịch vụ công khai cho hạ tầng trung gian, nhằm tối thiểu hóa chi phí vận hành ở giai đoạn nguyên mẫu:

    \begin{itemize}
        \item \textbf{Backend} được đóng gói bằng Docker và triển khai trên nền tảng đám mây (Render), chạy bằng máy chủ ASGI Uvicorn. Cấu hình được nạp qua biến môi trường (chuỗi kết nối MongoDB, khóa bí mật JWT, thông tin Telegram, tham số geofence mặc định).
        \item \textbf{Cơ sở dữ liệu} dùng MongoDB Atlas gói miễn phí, có sao lưu tự động.
        \item \textbf{Broker MQTT} dùng EMQX public broker (\texttt{broker.emqx.io}), không cần tự vận hành hạ tầng.
        \item \textbf{Thiết bị} nạp firmware qua PlatformIO và hoạt động độc lập ngoài hiện trường, kết nối Internet qua SIM dữ liệu (APN Viettel).
        \item \textbf{Ứng dụng di động} đóng gói thành tệp cài đặt (APK cho Android) để cài trực tiếp lên thiết bị phụ huynh.
    \end{itemize}

    Trong môi trường phát triển, toàn bộ hạ tầng phụ trợ (MongoDB và broker Mosquitto cục bộ) có thể khởi chạy nhanh bằng Docker Compose, giúp kiểm thử mà không phụ thuộc dịch vụ ngoài. Kết quả triển khai thử nghiệm cho thấy hệ thống vận hành ổn định với độ trễ cập nhật vị trí đầu cuối dưới 3 giây và gửi cảnh báo gần như tức thì trong điều kiện mạng di động thông thường.

    \end{document}
